\فصل{مقدمه}

\قسمت{بیان مسئله و ضرورت پژوهش}

مسئله زمان‌بندی کارگاهی\پاورقی{Job Shop Scheduling Problem (JSSP)}
یکی از مسائل بنیادی و پرکاربردِ بهینه‌سازی ترکیبیاتی\پاورقی{Combinatorial Optimization}
در حوزه‌های تحقیق در عملیات\پاورقی{Operations Research}
و مهندسی صنایع و رایانه است~\cite{pinedo2016scheduling}.
در این مسئله، مجموعه‌ای از کارها\پاورقی{Jobs} باید روی مجموعه‌ای از ماشین‌ها\پاورقی{Machines} پردازش شوند.
هر کار شامل دنباله‌ای مرتب از عملیات\پاورقی{Operations} است
و هر عملیات باید روی ماشین مشخصی و با زمان پردازش\پاورقی{Processing Time} معینی اجرا شود.
دو محدودیت اصلی بدین صورت است که
(۱) هر ماشین در هر لحظه حداکثر یک عملیات را پردازش می‌کند
و (۲) ترتیب عملیات هر کار باید رعایت شود (قید تقدم\پاورقی{Precedence Constraint}).
هدف کلاسیک این مسئله، کمینه‌سازی زمان تکمیل کل\پاورقی{Makespan}
(زمان اتمام آخرین عملیات) است.

مسئله زمان‌بندی کارگاهی از جمله مسائل ان‌پی-سخت\پاورقی{NP-hard} است~\cite{garey1976complexity}؛
ازاین‌رو یافتن پاسخ بهینه، حتی برای نمونه‌هایی با ابعاد متوسط، در زمان چندجمله‌ای امکان‌پذیر نیست.
در نتیجه، در ادبیات پژوهش از روش‌های ابتکاری\پاورقی{Heuristic}
و فراابتکاری\پاورقی{Metaheuristic} برای دستیابی به راه‌حل‌های نزدیک‌به‌بهینه استفاده می‌شود.
از جمله رویکردهای پرکاربرد می‌توان به الگوریتم ژنتیک\پاورقی{Genetic Algorithm}~\cite{goldberg1989genetic}،
تبرید شبیه‌سازی‌شده\پاورقی{Simulated Annealing}
و بهینه‌سازی ازدحام ذرات\پاورقی{Particle Swarm Optimization} اشاره کرد.

از سوی دیگر، در کاربردهای صنعتیِ نوین
بهینه‌سازی صرفِ زمان تکمیل کل کفایت نمی‌کند
و معیارهای دیگری نظیر مصرف انرژی\پاورقی{Energy Consumption}
و قابلیت اطمینان\پاورقی{Reliability} سامانه نیز اهمیت می‌یابند.
بنابراین، صورت‌بندی مسئله به‌صورت چندهدفه\پاورقی{Multi-Objective}
و به‌کارگیری الگوریتم‌های بهینه‌سازی چندهدفه، ضرورتی اجتناب‌ناپذیر است.

افزون بر این، با گسترش فناوری‌های
رایانش لبه‌ای\پاورقی{Edge Computing}~\cite{shi2016edge}
و رایانش مه\پاورقی{Fog Computing}~\cite{bonomi2012fog}،
زیرساخت‌های رایانشی ناهمگنِ چندلایه‌ای شکل گرفته‌اند
که گره‌های آن‌ها از نظر توان پردازشی، مصرف انرژی و نرخ خرابی
رفتارهای متفاوتی دارند.
در چنین بستری، تصمیم‌گیری صرفاً به تعیین ترتیب عملیات محدود نیست؛
بلکه نگاشت\پاورقی{Mapping} اجرای عملیات به گره‌های لبه، مه و ابر نیز
به‌عنوان بخشی از مسئله مطرح می‌شود.
این تلفیقِ «زمان‌بندی» و «تخصیص گره»، چالشی تازه ایجاد می‌کند
که در ادبیات کلاسیک زمان‌بندی کمتر به آن پرداخته شده است.

با توجه به مطالب فوق، \مهم{خلأ پژوهشی} را می‌توان چنین جمع‌بندی کرد:
بخش عمده‌ای از پژوهش‌های موجود، مسئله زمان‌بندی کارگاهی را به‌صورت تک‌هدفه مدل کرده
یا زیرساخت اجرای عملیات را همگن (یا مستقل از تصمیم زمان‌بندی) فرض می‌کنند.
در مقابل، پژوهش‌های اندکی به‌طور هم‌زمان
بهینه‌سازی سه‌هدفه را در بستر یک زیرساخت ناهمگنِ لبه-مه-ابر
در نظر گرفته‌اند.
ازاین‌رو، این پایان‌نامه با هدف پر کردن این خلأ،
به طراحی و ارزیابی یک الگوریتم ژنتیک چندهدفهٔ بهبودیافته
مبتنی بر \lr{NSGA-II}
برای حل مسئله زمان‌بندی کارگاهی در بستر لبه-مه-ابر می‌پردازد.


\قسمت{اهمیت و ضرورت پژوهش}

اهمیت پژوهش حاضر از جنبه‌های زیر قابل تبیین است:

\شروع{فقرات}
\فقره
\مهم{کاربردهای صنعتی گسترده}:
مسئله زمان‌بندی کارگاهی در سامانه‌های تولیدی،
خطوط مونتاژ و حتی مراکز داده\پاورقی{Data Centers} کاربرد دارد.
بهبود زمان‌بندی می‌تواند به کاهش هزینه‌های عملیاتی
و افزایش بهره‌وری منجر شود.
\فقره
\مهم{چالش محاسباتی}:
ان‌پی-سخت بودن مسئله سبب شده طراحی الگوریتم‌های کارا
برای آن همواره یک مسئله پژوهشی مهم باشد~\cite{garey1979computers}.
\فقره
\مهم{ضرورت بهینه‌سازی چندهدفه}:
در مسائل واقعی،
تصمیم‌گیرندگان با اهداف متعارض مواجه‌اند
و نیازمند مجموعه‌ای از راه‌حل‌های پارتو-بهینه\پاورقی{Pareto-optimal} هستند.
\فقره
\مهم{زیرساخت‌های رایانشی نوین}:
ظهور رایانش لبه‌ای و مه
ایجاب می‌کند که مدل‌های زمان‌بندی، ناهمگنی زیرساخت
و هزینه‌های ارتباطی میان لایه‌ها را نیز لحاظ کنند.
\پایان{فقرات}


\قسمت{اهداف پژوهش}

اهداف اصلی این پایان‌نامه عبارت‌اند از:

\شروع{شمارش}
\فقره
مدل‌سازی مسئله زمان‌بندی کارگاهی
در بستر زیرساخت رایانشی لبه-مه-ابر
با سه تابع هدف:
کمینه‌سازی زمان تکمیل کل،
کمینه‌سازی مصرف انرژی
و کمینه‌سازی عدم‌قابلیت‌اطمینان (معادلِ بیشینه‌سازی قابلیت اطمینان).
\فقره
طراحی و پیاده‌سازی یک الگوریتم ژنتیک چندهدفهٔ بهبودیافته
مبتنی بر \lr{NSGA-II}
مجهز به مقداردهی اولیه هوشمند\پاورقی{Smart Initialization}،
انتخاب تطبیقی عملگر\پاورقی{Adaptive Operator Selection}
و جست‌وجوی محلی\پاورقی{Local Search}.
\فقره
ارزیابی تجربی جامع الگوریتم پیشنهادی
روی ۱۵ نمونه معیار\پاورقی{Benchmark} استاندارد
و مقایسه آماری با روش‌های پایه.
\فقره
انجام مطالعه حذفی\پاورقی{Ablation Study}
برای تعیین سهم هر یک از مؤلفه‌های بهبودی
در عملکرد نهایی الگوریتم.
\پایان{شمارش}


\قسمت{مثال انگیزشی}

برای روشن‌تر شدن انگیزهٔ کاربردی پژوهش، سناریوی زیر را در نظر بگیرید:
چندین سفارش تولیدی به‌صورت هم‌زمان وارد سامانه می‌شوند و هر سفارش شامل
دنباله‌ای از عملیات وابسته است.
در چنین شرایطی، سامانه تصمیم‌گیری باید علاوه‌بر تعیین ترتیب اجرای عملیات روی ماشین‌ها،
تخصیص اجرای هر عملیات را نیز به گره‌های ناهمگنِ لبه، مه و ابر انجام دهد.
این سناریو در شکل~\رجوع{شکل:مثال-انگیزشی} به‌صورت شماتیک نمایش داده شده است.

\شروع{شکل}[H]
\centerimg{motivational_scenario.pdf}{14cm}
\شرح{مثال انگیزشی از زمان‌بندی کارگاهی چندهدفه در بستر لبه-مه-ابر و پیامدهای تصمیم‌گیری چندمعیاره}
\برچسب{شکل:مثال-انگیزشی}
\پایان{شکل}

همان‌گونه که در شکل مشاهده می‌شود، تصمیم‌های زمان‌بندی مستقیماً بر سه شاخص کلیدی اثر می‌گذارند:
زمان تکمیل کل، مصرف انرژی و قابلیت اطمینان.
اگر تصمیم‌گیری صرفاً بر مبنای کمینه‌سازی زمان تکمیل کل انجام شود،
در بسیاری از موارد بار پردازشی به گره‌های پرتوان‌تر منتقل می‌شود که می‌تواند
مصرف انرژی را افزایش دهد؛
از سوی دیگر، تأکید افراطی بر کاهش انرژی ممکن است
به افت عملکرد زمانی یا کاهش قابلیت اطمینان منجر شود.

ازاین‌رو، رویکرد چندهدفه مبتنی بر جبهه پارتو برای این دسته از مسائل ضروری است؛
زیرا به‌جای ارائهٔ یک راه‌حل منفرد،
مجموعه‌ای از راه‌حل‌های رقابتی و قابل انتخاب را
در اختیار تصمیم‌گیرنده قرار می‌دهد.
پژوهش حاضر با تمرکز بر همین نیاز عملی،
در پی ارائه روشی است که بتواند
موازنه‌ای مناسب میان معیارهای زمان، انرژی و قابلیت اطمینان برقرار کند.


\قسمت{نوآوری و سهم‌های پژوهش}

سهم‌های اصلی این پایان‌نامه به شرح زیر فهرست می‌شوند:

\شروع{شمارش}
\فقره
\مهم{مدل‌سازی سه‌هدفه در بستر لبه-مه-ابر}:
ارائه یک مدل یکپارچه که مسئله زمان‌بندی کارگاهی را
با شبیه‌سازی زیرساخت ناهمگن ترکیب می‌کند
و سه هدف زمان تکمیل کل، مصرف انرژی و قابلیت اطمینان
را به‌صورت هم‌زمان بهینه می‌کند.
\فقره
\مهم{الگوریتم ژنتیک بهبودیافته}:
ترکیب سه تکنیک بهبودی
(مقداردهی اولیه هوشمند مبتنی بر قواعد ابتکاری
\lr{SPT}\پاورقی{Shortest Processing Time}
و \lr{MWR}\پاورقی{Most Work Remaining}،
انتخاب تطبیقی عملگر مبتنی بر \lr{UCB1}
و جست‌وجوی محلی)
در چارچوب \lr{NSGA-II}.
\فقره
\مهم{ارزیابی تجربی جامع}:
آزمایش روی ۱۵ نمونه معیار استاندارد
(شامل خانواده‌های \lr{FT}، \lr{LA}، \lr{ORB}، \lr{ABZ}، \lr{TA}، \lr{SWV} و \lr{YN})
در قالب ۱۰ اجرای مستقل،
همراه با آزمون‌های آماری ناپارامتری
(فریدمن و ویلکاکسون با تصحیح هولم)
و مطالعه حذفی.
\فقره
\مهم{قابلیت بازتولید نتایج}:
کدها، داده‌ها و پیکربندی‌های لازم
به‌گونه‌ای سامان‌دهی و مستندسازی شده‌اند
که بازتولید آزمایش‌ها با بذرهای تصادفی ثابت امکان‌پذیر باشد.
\پایان{شمارش}


\قسمت{ساختار پایان‌نامه}

ادامه این پایان‌نامه در چهار فصل دیگر سازمان‌دهی شده است.
در \مهم{فصل دوم} مفاهیم پایه و مبانی نظری شامل
مدل ریاضی رسمی مسئله زمان‌بندی کارگاهی،
اصول بهینه‌سازی چندهدفه و الگوریتم \lr{NSGA-II}
و مدل زیرساخت رایانشی لبه-مه-ابر ارائه می‌شوند.
\مهم{فصل سوم} به مرور ادبیات و پیشینه پژوهش اختصاص دارد
و شامل بررسی رویکردهای دقیق، ابتکاری و فراابتکاری
برای حل مسئله زمان‌بندی کارگاهی،
مرور پژوهش‌های فارسی مرتبط
و بررسی کارهای مرتبط با بهینه‌سازی چندهدفه و زمان‌بندی در بستر لبه-مه-ابر است.
در \مهم{فصل چهارم} روش پیشنهادی این پایان‌نامه
شامل نمایش کروموزوم، عملگرهای ژنتیکی و مؤلفه‌های بهبودی تشریح می‌شود
و سپس نتایج آزمایش‌ها روی نمونه‌های معیار ارائه و تحلیل می‌شوند.
در نهایت، \مهم{فصل پنجم} به جمع‌بندی نتایج،
بحث درباره یافته‌ها، بیان محدودیت‌ها
و ارائه پیشنهادهایی برای پژوهش‌های آتی اختصاص دارد.


\قسمت{جمع‌بندی فصل}

در این فصل، مسئله زمان‌بندی کارگاهی و ضرورت پرداختن به آن در بسترهای
رایانشی نوین تبیین شد.
همچنین خلأ پژوهشی، اهداف، سهم‌های اصلی و ساختار کلی پایان‌نامه ارائه شد.
بر این اساس، فصل بعد به بیان مبانی نظری و مدل‌های رسمی مورد استفاده در پژوهش اختصاص دارد.
